% Inclusion-ready proof draft. Requires the parent's cq macros/theorem environments.
% This file was written independently; the parent decides the insertion point.
\subsection{The exact Schwartz--Mellin image}
\label{sec:cq-Schwartz-Mellin-image}
Put $\theta=x\partial_x$. For real $a\leq b$, retain the finite polynomial
\[
 P_{a,b}(z)=\prod_{\substack{k\geq1\\a-1\leq-2k\leq b+1}}(z+2k),
\]
with empty product equal to $1$. Define $\mathfrak M_0^{\rm ev}$
to be the vector space of meromorphic functions $M$ on $\cqC$
whose only possible poles are simple poles at $-2k$, $k\geq1$,
which satisfy $M(1)=0$, and for which all the seminorms
\begin{equation}\label{eq:cq-Mellin-strip-seminorm}
 q_{a,b,N}(M)=
 \sup_{\substack{a\leq\sigma\leq b\\v\in\mathbb R}}
 (1+|v|)^N
 \bigl|P_{a,b}(\sigma+iv)M(\sigma+iv)\bigr|,
 \qquad N=0,1,2,\ldots,
\end{equation}
are finite. At a pole of $M$, the product in this formula has its
removable value. The margin $1$ in $P_{a,b}$ is retained in this
definition. It ensures that all poles in a neighborhood of the
closed strip have been removed. These seminorm conditions are
equivalent to rapid decrease of $M$ on every finite vertical strip
for $|\operatorname{Im}z|\geq1$: on that part of the strip each
factor $|z+2k|$ is bounded above by a constant times $1+|v|$
and is at least $|v|$; on the remaining compact rectangle,
$P_{a,b}M$ is holomorphic in a neighborhood and is bounded.
In particular no bound uniform in the pole index $k$ has been
inserted into the definition.
One can equivalently use the residue coordinates together with
the explicit away-from-poles strip seminorms
\[
 r_{a,b,N,\delta}(M)=
 \sup_{\substack{a\leq\operatorname{Re}z\leq b\\
       \operatorname{dist}(z,\{-2,-4,\ldots\})\geq\delta}}
   (1+|\operatorname{Im}z|)^N|M(z)|,
 \qquad 0<\delta\leq\tfrac12.
\]
Here is the equivalence as a statement about topologies, not only
about membership. On the domain of this last supremum every
factor of $P_{a,b}$ has modulus at least $\delta$, so
$r_{a,b,N,\delta}\leq\delta^{-\deg P_{a,b}}q_{a,b,N}$.
Residues are controlled by $q$ as shown explicitly below.
Conversely, outside the radius-$\delta$ disks about the poles,
$|P_{a,b}(z)|\leq C_{a,b}(1+|\operatorname{Im}z|)^{\deg P_{a,b}}$
on the strip, so the corresponding part of $q$ is bounded by
an $r$ seminorm with this degree added to its exponent.
For a pole whose radius-$\delta$ disk meets the strip, that pole
is among the factors of $P_{a,b}$. The function $P_{a,b}M$
is holomorphic throughout the disk; its boundary has distance
at least $\delta$ from every pole, since distinct pole locations
are separated by $2$. That boundary lies in the enlarged strip
$a-1\leq\operatorname{Re}z\leq b+1$. The maximum-modulus
principle bounds the product inside the disk by its boundary
values, and $(1+|\operatorname{Im}z|)^N\leq(1+\delta)^N$
there. An $r$ seminorm on that enlarged strip therefore bounds
the remaining part of $q$. There are finitely many relevant disks.
These estimates prove equality of the two locally convex
topologies; separate residue coordinates may be included without
changing them.

The map
\begin{equation}\label{eq:cq-Mellin-isomorphism}
 \mathcal M:\mathcal S_0^{\rm ev}\longrightarrow
                 \mathfrak M_0^{\rm ev},\qquad
 (\mathcal M\psi)(z)=\int_0^\infty\psi(x)x^z\,\frac{dx}{x}
 \quad(\operatorname{Re}z>-2),
\end{equation}
is a linear bijection. Its inverse, for every fixed real $c>1$, is
\begin{equation}\label{eq:cq-Mellin-reconstruction}
 \psi(x)=\frac1{2\pi}\int_{\mathbb R}
              M(c+iv)x^{-c-iv}\,dv\quad(x>0),
 \qquad \psi(0)=0,\qquad\psi(-x)=\psi(x).
\end{equation}
For every $k\geq1$ its precise residue is
\begin{equation}\label{eq:cq-Mellin-Taylor-residue}
 \operatorname*{Res}_{z=-2k}M(z)
       =\frac{\psi^{(2k)}(0)}{(2k)!}.
\end{equation}
Both maps are continuous when the source has its ordinary Schwartz
seminorms and the target has \eqref{eq:cq-Mellin-strip-seminorm}.

Here is a proof including the strip and reconstruction estimates.
Fix a smooth function $\chi$ on $[0,\infty)$ which equals $1$
on $[0,1]$ and $0$ on $[2,\infty)$. Such a function is obtained,
for example, by setting $\eta(y)=0$ for $y\leq0$ and
$\eta(y)=e^{-1/y}$ for $y>0$, and taking
$\chi(x)=\eta(2-x)/(\eta(2-x)+\eta(x-1))$. Its denominator is
positive at every $x\geq0$. Let
$a_k=\psi^{(2k)}(0)/(2k)!$. For an integer $K\geq1$ write
\[
 r_K(x)=\psi(x)-\chi(x)\sum_{k=1}^K a_kx^{2k}.
\]
Evenness and $\psi(0)=0$ show that $r_K$ vanishes to order
at least $2K+2$ at zero. Taylor's integral remainder, applied
also after differentiating, proves
\[
 |\theta^Jr_K(x)|\leq C_{J,K}(\psi)x^{2K+2}
 \quad(0<x\leq1),\qquad J\geq0,
\]
where $C_{J,K}(\psi)$ is bounded by a constant times finitely
many Schwartz seminorms of $\psi$. To see explicitly that all
$J$ are covered, expand
$\theta^J=\sum_{\ell=0}^J S(J,\ell)x^\ell\partial_x^\ell$,
where $S(0,0)=1$ and
$S(J+1,\ell)=\ell S(J,\ell)+S(J,\ell-1)$. For
$\ell\leq2K+2$, Taylor's remainder gives
\[
 |\partial_x^\ell r_K(x)|
 \leq\frac{x^{2K+2-\ell}}{(2K+2-\ell)!}
              \sup_{0\leq y\leq1}|\psi^{(2K+2)}(y)|.
\]
For $\ell>2K+2$, the differentiated polynomial vanishes on
this interval and $x^\ell\leq x^{2K+2}$, which gives the
claimed estimate using $\sup_{[0,1]}|\psi^{(\ell)}|$.
At infinity $r_K=\psi$, so each $\theta^Jr_K$ decreases
faster than every inverse power.

The cutoff Mellin function
\[
 C(w)=\int_0^\infty\chi(x)x^w\,\frac{dx}{x}
       \quad(\operatorname{Re}w>0)
\]
has the meromorphic continuation
\begin{equation}\label{eq:cq-Mellin-cutoff}
 C(w)=-\frac{L(w)}{w},\qquad
 L(w)=\int_1^2\chi'(x)x^w\,dx.
\end{equation}
The entire function $L$ satisfies $L(0)=\chi(2)-\chi(1)=-1$,
so the residue of $C$ at $0$ is $1$. Repeated integration by
parts with $\theta$ in the expression
$L(w)=\int_1^2(x\chi'(x))x^w\,dx/x$ gives, on every finite
vertical strip, rapid decrease for $|\operatorname{Im}w|\geq1$.
Indeed $x\chi'(x)$ is smooth with support in $[1,2]$ and all
boundary derivatives vanish, and for every integer $J\geq0$
\[
 w^JL(w)=(-1)^J\int_1^2
       \theta^J(x\chi'(x))x^w\,\frac{dx}{x}.
\]
Thus $C$ has the same rapid decrease away from its sole pole.

For any strip $a\leq\operatorname{Re}z\leq b$, choose $K$
with $a+2K+2>0$. On the larger half-plane
$\operatorname{Re}z>-2K-2$, the formula
\begin{equation}\label{eq:cq-Mellin-meromorphic-extension}
 M(z)=\int_0^\infty r_K(x)x^z\,\frac{dx}{x}
                     +\sum_{k=1}^K a_k C(z+2k)
\end{equation}
is meromorphic and agrees with the original Mellin integral where
$\operatorname{Re}z>-2$. The expressions for successive $K$
therefore agree by the identity theorem and define one meromorphic
function on $\cqC$. Formula \eqref{eq:cq-Mellin-cutoff} proves
that its only possible poles and residues are exactly
\eqref{eq:cq-Mellin-Taylor-residue}.
For every $J\geq0$, integration by parts gives
\[
 z^J\int_0^\infty r_K(x)x^z\,\frac{dx}{x}
       =(-1)^J\int_0^\infty\theta^Jr_K(x)x^z\,\frac{dx}{x}.
\]
All boundary terms vanish: at zero they are bounded by a constant
times $x^{a+2K+2}$, and at infinity the Schwartz exponent can
be chosen larger than $b+J+1$. The absolute value of the final
integral, uniformly in this strip, is at most
\begin{equation}\label{eq:cq-Mellin-forward-bound}
 B_{K,J;a,b}(\psi)=
 \int_0^1|\theta^Jr_K(x)|x^{a-1}\,dx
 +\int_1^\infty|\theta^Jr_K(x)|x^{b-1}\,dx<\infty.
\end{equation}
For $|\operatorname{Im}z|\geq1$, divide by $|z|^J\geq
|\operatorname{Im}z|^J$. Taking $J$ arbitrarily large proves
all the required strip estimates, including the polynomial
$P_{a,b}$. At bounded imaginary part, multiply
\eqref{eq:cq-Mellin-meromorphic-extension} by $P_{a,b}$;
every pole meeting the strip cancels and the resulting expression
is bounded on the compact rectangle. The constants in these
bounds involve only finitely many Schwartz seminorms:
the first integral in \eqref{eq:cq-Mellin-forward-bound} is
at most $C_{J,K}(\psi)/(a+2K+2)$; on $[1,2]$ the cutoff
has fixed bounded derivatives; on $[2,\infty)$ the expansion
of $\theta^J$ above uses finitely many bounds
$(1+x)^A|\psi^{(\ell)}(x)|$ with $A>b+J+1$.
This proves continuity of the forward map.
Finally $M(1)=\int_0^\infty\psi(x)\,dx
=\widehat\psi(0)/2=0$, retaining the factor $2$ from evenness.

Conversely, let $M\in\mathfrak M_0^{\rm ev}$ and define
\eqref{eq:cq-Mellin-reconstruction}. The integral and its
ordinary $x$ derivatives converge uniformly on each compact
subinterval of $(0,\infty)$. More precisely, with the rising
factorial $(z)_0=1$ and
$(z)_j=z(z+1)\cdots(z+j-1)$ for $j\geq1$, they are
\begin{equation}\label{eq:cq-Mellin-differentiated-inverse}
 \psi^{(j)}(x)=\frac{(-1)^j}{2\pi}
       \int_{\mathbb R}(c+iv)_jM(c+iv)x^{-c-iv-j}\,dv.
\end{equation}
The factor $(c+iv)_j$ grows at most polynomially in $v$,
so a strip estimate with exponent greater than $j+1$ supplies
an integrable dominating function. This proves smoothness for $x>0$.

The inverse integral is independent of the positive choice of
$c$. In fact integrate $M(z)x^{-z}$ around the rectangle
between two positive vertical lines and heights $\pm T$.
It has no poles there. On each horizontal side, its absolute
value is at most a constant, depending on the fixed $x$ and
the bounded real interval, times $(1+T)^{-N}$ for arbitrary
$N$. Thus the horizontal contributions tend to zero, and the
two upward vertical integrals agree. In particular for every
real $b>0$ the differentiated formula holds with $c$ replaced
by $b$. If
\begin{equation}\label{eq:cq-Mellin-line-integral}
 I_{d,j}(M)=\frac1{2\pi}\int_{\mathbb R}
                |(d+iv)_jM(d+iv)|\,dv
 \quad\bigl(d\notin\{-2,-4,-6,\ldots\}\bigr),
\end{equation}
then $I_{d,j}(M)$ is finite and
\begin{equation}\label{eq:cq-Mellin-infinity-bound}
 |\psi^{(j)}(x)|\leq I_{b,j}(M)x^{-b-j}
 \quad(x>0,\ b>0).
\end{equation}
This gives every required decay estimate at infinity.

For the behavior at zero fix $K\geq1$ and take the actual
left vertical line $d=-2K-1$. The rectangle between $d$ and
$c$ has precisely the possible poles $-2,-4,\ldots,-2K$.
Traverse its right edge upward and its left edge downward;
this is the positive orientation. The residue theorem and
the same horizontal estimates give
\begin{equation}\label{eq:cq-Mellin-contour-shift}
 \begin{aligned}
 \psi(x)&=\sum_{k=1}^K a_kx^{2k}+R_K(x),
 &a_k&=\operatorname*{Res}_{z=-2k}M(z),\\
 R_K(x)&=\frac1{2\pi}\int_{\mathbb R}
                    M(-2K-1+iv)x^{2K+1-iv}\,dv,\\
 |R_K^{(j)}(x)|&\leq I_{-2K-1,j}(M)x^{2K+1-j}.
 \end{aligned}
\end{equation}
The sign of each residue is positive. The last estimate follows
by differentiating the remainder integral with the exact factor
$(-1)^j(-2K-1+iv)_j$; this factor is included in
\eqref{eq:cq-Mellin-line-integral}.
For any given nonnegative integer $j$, choose $K$ so that
$2K+1>j$. Formula \eqref{eq:cq-Mellin-contour-shift} then
proves that all derivatives through order $j$ have limits at zero
given by the displayed even polynomial. These limits define a
$C^j$ extension: inductively, if $\psi^{(\ell)}$ and
$\psi^{(\ell+1)}$ have these limits, integrate
$\psi^{(\ell+1)}$ from a positive $\varepsilon$ to $x$ and
let $\varepsilon\downarrow0$; the resulting integral identity
proves differentiability at zero with the asserted derivative.
Because $j$ was arbitrary, the extension is smooth, has value
zero, all its odd derivatives vanish at zero, and its $2k$th
derivative is $(2k)!a_k$. Reflecting it evenly across zero
preserves smoothness: on the negative half-line its $j$th
derivative is $(-1)^j\psi^{(j)}(-x)$, and the one-sided
values coincide because the odd jets vanish. Together with
\eqref{eq:cq-Mellin-infinity-bound}, this proves
$\psi\in\mathcal S(\mathbb R)$ and the residue assertion.

For completeness the inversion just used also proves that the
Mellin transform of this reconstructed function is the prescribed
$M$. On a line $c>1$ put
$h(u)=e^{cu}\psi(e^u)$ and $m(v)=M(c+iv)$. The formula says
$h(u)=(2\pi)^{-1}\int m(v)e^{-ivu}\,dv$.
Here $m$ is continuous and integrable by the strip estimates,
and $h$ is integrable by the bounds at both ends just proved.
For $\varepsilon>0$, absolute Fubini and the Gaussian integral give
\[
 \int_{\mathbb R}h(u)e^{ivu}e^{-\varepsilon u^2/2}\,du
 =\int_{\mathbb R}m(w)
       \frac{e^{-(v-w)^2/(2\varepsilon)}}{\sqrt{2\pi\varepsilon}}
       \,dw.
\]
The right side tends to $m(v)$: split the integral at
$|v-w|=\delta$, use continuity on the inner interval, and use
the Gaussian tail and boundedness of $m$ on its complement.
The left side tends by dominated convergence to
$\int h(u)e^{ivu}\,du$. Substitution $x=e^u$ proves
$\mathcal M\psi(c+iv)=M(c+iv)$. Both sides are holomorphic
on $\operatorname{Re}z>-2$, so the identity theorem proves
their equality there and hence equality of their meromorphic
continuations. In particular
$\widehat\psi(0)=2\mathcal M\psi(1)=2M(1)=0$.
Reconstruction is injective because its output has Mellin transform
equal to its input. The forward map is injective as well, with
the following direct uniqueness proof. If the Mellin transform of
$\delta\in\mathcal S_0^{\rm ev}$ vanishes, the continuous
integrable function $g(u)=e^{cu}\delta(e^u)$ has Fourier transform
$\int g(u)e^{ivu}\,du=0$. Put
$G_\varepsilon(u)=(2\pi\varepsilon)^{-1/2}
e^{-u^2/(2\varepsilon)}$. Absolute Fubini applied to the Gaussian
integral gives
\[
 (g*G_\varepsilon)(u)=\frac1{2\pi}
 \int_{\mathbb R}\left(\int_{\mathbb R}g(w)e^{ivw}\,dw\right)
                   e^{-\varepsilon v^2/2}e^{-ivu}\,dv=0.
\]
As $\varepsilon\downarrow0$, the Gaussian convolution tends
to $g(u)$ at every $u$ by continuity, the unit total mass of
$G_\varepsilon$, and its vanishing tail mass. Thus $g=0$,
$\delta=0$ on the positive half-line, and evenness gives
$\delta=0$ everywhere. This proves that the two displayed maps
are mutual inverses, not only that reconstructed inputs have the
prescribed transform.

The inverse continuity can be read directly from the estimates.
For nonnegative integers $A,j$, choose $b=A+1$ and $K\geq1$
with $2K+1>j$. On $|x|\geq1$ use
\eqref{eq:cq-Mellin-infinity-bound}; on $|x|\leq1$ use
\eqref{eq:cq-Mellin-contour-shift}. They give
\begin{equation}\label{eq:cq-Mellin-inverse-seminorm}
 \sup_{x\in\mathbb R}(1+|x|)^A|\psi^{(j)}(x)|
 \leq2^A\left(
 I_{b,j}(M)+I_{-2K-1,j}(M)
 +\sum_{\substack{1\leq k\leq K\\2k\geq j}}
           \frac{(2k)!}{(2k-j)!}|a_k|\right).
\end{equation}
Each line integral on the right is bounded by a constant times
a strip seminorm with exponent greater than $j+1$ on a strip
containing its line and no pole on the line. Every one of the
finitely many residues is controlled by a strip seminorm too:
if $-2k$ belongs to $[a,b]$, then
\[
 |a_k|=
 \frac{|(P_{a,b}M)(-2k)|}
 {\left|\displaystyle\prod_{\substack{\ell\geq1,\ \ell\ne k\\
                         a-1\leq-2\ell\leq b+1}}(2\ell-2k)\right|}.
\]
Its denominator is a finite product of nonzero numbers. This
proves continuity without a uniform assumption on the entire
residue sequence and completes the characterization.

Apply this result to the exact arithmetic trace map already proved
in \eqref{eq:cq-Schwartz-trace}. For $F\in\cqE$ define the
meromorphic function
\[
 M_F(z)=\frac{F(i(z-\tfrac12))}{\zeta(z)}.
\]
Its definition uses meromorphic division, including the removable
values determined by any canceled zeros or pole. Then the precise
intersection of the full source trace image with its original
order-at-most-one entire-function set is
\begin{equation}\label{eq:cq-exact-Schwartz-trace-image}
 \begin{aligned}
 \mathcal A_{\rm Sch}
 &:=\{F_\psi:\psi\in\mathcal S_0^{\rm ev}\}\cap\cqE\\
 &=\{F\in\cqE:M_F\in\mathfrak M_0^{\rm ev}\}\\
 &=\left\{\Xi h:h\in\cqE,\quad
     A(z)h(i(z-\tfrac12))\in\mathfrak M_0^{\rm ev}\right\},\\
 A(z)&=\tfrac12z(z-1)\pi^{-z/2}\Gamma(z/2).
 \end{aligned}
\end{equation}
For the first equality after the definition, a trace has the
factorization and Mellin properties already proved. Conversely,
if $M_F\in\mathfrak M_0^{\rm ev}$, reconstruct the unique
$\psi\in\mathcal S_0^{\rm ev}$. In $\operatorname{Re}z>1$
the trace factorization gives
$F_\psi(i(z-\tfrac12))=\zeta(z)M_F(z)
=F(i(z-\tfrac12))$. Both $F_\psi$ and $F$ are entire;
their equality on this open set proves equality everywhere.
This argument retains the original $s$ coordinate and does not
assert that arbitrary Schwartz traces belong to $\cqE$.

For the last equality in
\eqref{eq:cq-exact-Schwartz-trace-image}, the previous
nontrivial-zero jet calculation and the proved entire division
theorem give $F=\Xi h$ with $h\in\cqE$ for every trace in
this intersection. Conversely the retained identity
$\Xi(i(z-\tfrac12))=\xi(1-z)=\xi(z)$ gives exactly
\[
 M_{\Xi h}(z)=
      \tfrac12z(z-1)\pi^{-z/2}\Gamma(z/2)
                      h(i(z-\tfrac12)),
\]
including its meromorphic continuation and all its local units.
The factor $A$ has only simple poles at $-2k$, $k\geq1$:
the gamma poles at $z=0$ and the factor $z$ cancel, the
remaining gamma poles occur at the displayed negative even
integers, and the factor $z-1$ gives $A(1)=0$.
Thus for an arbitrary $h\in\cqE$ the pole locations,
orders and the equality $M_{\Xi h}(1)=0$ are automatic;
the exact additional requirement is the family of strip bounds
in \eqref{eq:cq-Mellin-strip-seminorm}, with the actual factor
$A$ retained. No topology on the unnamed source closure is
identified by this characterization.

This criterion also retains every Taylor coefficient of its
arithmetic preimage. The trivial zero of $\zeta$ at $-2k$ is
simple, as follows directly from the retained completion.
The gamma recurrence gives a simple pole of $\Gamma(z/2)$
there with nonzero residue $2(-1)^k/k!$, while the other factors
in $A(z)$ are nonzero there. On the other hand
$\xi(-2k)=\xi(1+2k)>0$: all factors in the completion at $1+2k$
are positive, including the convergent positive Dirichlet series
for $\zeta(1+2k)$. Thus $\xi=A\zeta$ forces a zero of exactly
order one at $-2k$. Meromorphic division now gives
\begin{equation}\label{eq:cq-Schwartz-arithmetic-Taylor-data}
 \begin{aligned}
 \frac{\psi^{(2k)}(0)}{(2k)!}
 &=\frac{F(i(-2k-\tfrac12))}{\zeta'(-2k)}\\
 &=\frac{2k(2k+1)(-1)^k\pi^k}{k!}
                     h\bigl(-i(2k+\tfrac12)\bigr)
 \qquad(F=\Xi h\in\mathcal A_{\rm Sch}).
 \end{aligned}
\end{equation}
In the first line there is no extra factor $i$: the residue is
taken in the actual coordinate $z$ and the denominator has leading
term $\zeta'(-2k)(z+2k)$. For the second line the gamma
recurrence gives
$\operatorname*{Res}_{z=-2k}\Gamma(z/2)=2(-1)^k/k!$.
Indeed $\Gamma(w)=\Gamma(w+k+1)/
[w(w+1)\cdots(w+k)]$, whose numerator is $1$ at $w=-k$
and whose denominator has derivative $(-1)^kk!$ there;
substitution $w=z/2$ supplies the factor $2$. Multiplying
this exact residue by
$(-2k)(-2k-1)\pi^k/2$ and the value of the entire $h$
gives the asserted coefficient. Thus the inclusion criterion
also reconstructs the full even jet at the origin.

The inclusion $\mathcal A_{\rm Sch}\hookrightarrow\Xi\cqE$
is the actual identity on entire functions, with exact image
specified by \eqref{eq:cq-exact-Schwartz-trace-image}. It contains
$\cqC[s]\Xi$ because the retained finite Hermite trace families
span precisely that polynomial image. Consequently its closures
in each of the explicitly proved analytic topologies satisfy
\begin{equation}\label{eq:cq-Schwartz-image-analytic-closures}
 \overline{\mathcal A_{\rm Sch}}^{\,\gamma}
 =\overline{\mathcal A_{\rm Sch}}^{\,\kappa}
 =\Xi\cqE.
\end{equation}
Indeed the left image lies between $\cqC[s]\Xi$ and the
closed subspace $\Xi\cqE$, and the closures of the smaller
polynomial subspace were proved to equal that subspace.
This constructs the exact inclusion, its image criterion and both
analytic closures; it inserts no equality involving the unnamed
topology $\tau$.
